% !TEX program = xelatex
\documentclass[aspectratio=169, 11pt]{beamer}

% ── 字体 & 中文 ──────────────────────────────────────
\usepackage{fontspec}
\usepackage{xeCJK}
\setCJKmainfont{Noto Serif CJK SC}
\setCJKsansfont{Noto Sans CJK SC}
\setCJKmonofont{Noto Sans Mono CJK SC}

% ── 数学 ─────────────────────────────────────────────
\usepackage{amsmath, amssymb, bm}

% ── 图形 ─────────────────────────────────────────────
\usepackage{tikz}
\usetikzlibrary{arrows.meta, positioning, shapes.geometric, calc,
                decorations.pathreplacing, fit, backgrounds, shadows}

% ── 表格 ─────────────────────────────────────────────
\usepackage{booktabs}
\usepackage{colortbl}
\usepackage{array}

% ── Beamer 主题 ──────────────────────────────────────
\usetheme{default}
\usecolortheme{default}
\setbeamertemplate{navigation symbols}{}
\setbeamertemplate{footline}[frame number]

% ── 颜色定义 ─────────────────────────────────────────
\definecolor{vaecol}{HTML}{2980B9}
\definecolor{vqcol}{HTML}{C0392B}
\definecolor{greencol}{HTML}{27AE60}
\definecolor{orangecol}{HTML}{E67E22}
\definecolor{purplecol}{HTML}{8E44AD}
\definecolor{darkcol}{HTML}{2C3E50}
\definecolor{lightbg}{HTML}{ECF0F1}
\definecolor{lightblue}{HTML}{D6EAF8}
\definecolor{lightred}{HTML}{FADBD8}
\definecolor{titlebg}{HTML}{1A1A2E}

\setbeamercolor{frametitle}{fg=darkcol, bg=white}
\setbeamercolor{title}{fg=white, bg=titlebg}
\setbeamercolor{normal text}{fg=darkcol}
\setbeamercolor{itemize item}{fg=vaecol}
\setbeamercolor{itemize subitem}{fg=orangecol}
\setbeamerfont{frametitle}{size=\Large, series=\bfseries}

% ── 自定义框 ─────────────────────────────────────────
\newcommand{\vaebox}[1]{%
  \begin{tikzpicture}[baseline=(X.base)]
    \node[fill=lightblue, rounded corners=3pt, inner sep=5pt] (X) {#1};
  \end{tikzpicture}}
\newcommand{\vqbox}[1]{%
  \begin{tikzpicture}[baseline=(X.base)]
    \node[fill=lightred, rounded corners=3pt, inner sep=5pt] (X) {#1};
  \end{tikzpicture}}

% ── TikZ 样式 ────────────────────────────────────────
\tikzset{
  block/.style={rectangle, rounded corners=5pt, minimum height=1cm,
                minimum width=2cm, text=white, font=\bfseries,
                drop shadow={shadow xshift=1pt, shadow yshift=-1pt, opacity=0.3}},
  vaeb/.style={block, fill=vaecol},
  vqb/.style={block, fill=vqcol},
  decb/.style={block, fill=darkcol},
  genb/.style={block, fill=greencol},
  orb/.style={block, fill=orangecol},
  purb/.style={block, fill=purplecol},
  arr/.style={-{Stealth[length=6pt]}, thick, color=#1},
  arr/.default={darkcol!60},
  dasharr/.style={-{Stealth[length=6pt]}, thick, dashed, color=#1},
  dasharr/.default={greencol},
  labelnode/.style={font=\small, text=darkcol},
  infobox/.style={rectangle, rounded corners=4pt, draw=#1, fill=#1!8,
                  text width=5.5cm, inner sep=6pt, font=\small},
}

% ═════════════════════════════════════════════════════
\begin{document}

% ── Slide 1: 封面 ───────────────────────────────────
{
\setbeamercolor{background canvas}{bg=titlebg}
\begin{frame}[plain]
\begin{tikzpicture}[remember picture, overlay]
  % 装饰圆
  \foreach \i in {1,...,12}{
    \pgfmathsetmacro{\cx}{rand*7}
    \pgfmathsetmacro{\cy}{rand*4}
    \pgfmathsetmacro{\r}{0.1+rand*0.3}
    \fill[vaecol, opacity=0.08] (\cx, \cy) circle (\r);
  }
  \foreach \i in {1,...,10}{
    \pgfmathsetmacro{\cx}{rand*7}
    \pgfmathsetmacro{\cy}{rand*4}
    \pgfmathsetmacro{\s}{0.1+rand*0.25}
    \fill[vqcol, opacity=0.08, rounded corners=2pt]
      (\cx-\s, \cy-\s) rectangle (\cx+\s, \cy+\s);
  }
\end{tikzpicture}
\vspace{1.5cm}
\begin{center}
  {\Huge\bfseries\color{white} VQ-VAE \;\;vs\;\; VAE}\\[12pt]
  {\large\color{white!70} Vector Quantized VAE\;\; vs\;\; Variational Autoencoder}\\[20pt]
  {\color{vqcol}\rule{6cm}{1.5pt}}\\[14pt]
  {\large\color{white!80}\itshape --- 从连续到离散的潜在空间之旅 ---}\\[24pt]
  {\color{white!60} Deep Learning $\cdot$ Spring 2026}
\end{center}
\end{frame}
}

% ── Slide 2: 总览表格 ───────────────────────────────
\begin{frame}{一张表看懂 VAE vs VQ-VAE}
\vspace{-2pt}
\centering
\renewcommand{\arraystretch}{1.45}
\begin{tabular}{>{\bfseries\color{darkcol}}l
                >{\color{vaecol}}c
                >{\color{vqcol}}c}
\toprule
\rowcolor{lightbg}
\textcolor{darkcol}{对比维度}
  & \textcolor{vaecol}{\textbf{VAE}}
  & \textcolor{vqcol}{\textbf{VQ-VAE}} \\
\midrule
潜在空间
  & 连续 (Continuous)
  & 离散 (Discrete) \\
\rowcolor{lightbg}
潜在变量
  & $z \sim \mathcal{N}(\mu, \sigma^2)$
  & $z_q = e_k$ \;(码本向量) \\
先验分布
  & $\mathcal{N}(0, I)$
  & Uniform Categorical \\
\rowcolor{lightbg}
梯度传播
  & 重参数化技巧
  & Straight-Through Est. \\
正则项
  & KL 散度
  & Commitment Loss \\
\rowcolor{lightbg}
生成质量
  & 偏模糊
  & 更清晰 \\
生成方式
  & 直接采样先验
  & 需额外先验模型 \\
\bottomrule
\end{tabular}
\end{frame}

% ── Slide 3: VAE 架构 ───────────────────────────────
\begin{frame}{VAE 架构}
\centering
\begin{tikzpicture}[node distance=1.2cm, scale=0.92, every node/.style={scale=0.92}]
  % 节点
  \node[vaeb, minimum width=1.8cm] (input) {Input $x$};
  \node[vaeb, right=1.5cm of input, minimum width=2.2cm] (enc) {Encoder $q_\phi$};
  \node[vaeb, above right=0.3cm and 1.2cm of enc, minimum width=1.4cm, fill=vaecol!80] (mu) {$\mu$};
  \node[vaeb, below right=0.3cm and 1.2cm of enc, minimum width=1.4cm, fill=vaecol!80] (sig) {$\sigma$};
  \node[genb, right=2.8cm of enc, minimum width=2.4cm, minimum height=1.2cm] (reparam)
    {\begin{tabular}{c}\small 重参数化\\ $z{=}\mu{+}\sigma{\cdot}\varepsilon$\end{tabular}};
  \node[decb, right=1.5cm of reparam, minimum width=2.2cm] (dec) {Decoder $p_\theta$};

  % 箭头
  \draw[arr] (input) -- (enc);
  \draw[arr=vaecol] (enc.east) ++(0,0.15) -- (mu.west);
  \draw[arr=vaecol] (enc.east) ++(0,-0.15) -- (sig.west);
  \draw[arr=greencol] (mu.east) -- ++(0.4,0) |- ([yshift=3pt]reparam.west);
  \draw[arr=greencol] (sig.east) -- ++(0.4,0) |- ([yshift=-3pt]reparam.west);
  \draw[arr] (reparam) -- (dec);

  % epsilon
  \node[above=0.5cm of reparam, font=\itshape\color{greencol}] (eps)
    {$\varepsilon \sim \mathcal{N}(0, I)$};
  \draw[arr=greencol] (eps) -- (reparam);

  % 潜在空间标注
  \node[below=1.0cm of reparam, fill=lightblue, rounded corners=4pt,
        inner sep=5pt, font=\bfseries\color{vaecol}] (zlabel)
    {连续潜在空间 $z \in \mathbb{R}^d$};

  % 损失函数
  \node[below=2.0cm of enc, xshift=2.5cm,
        draw=vaecol, rounded corners=5pt, fill=vaecol!5,
        inner sep=8pt, font=\normalsize, text=darkcol] (loss)
    {$\displaystyle \mathcal{L}_{\text{VAE}}
     = \underbrace{\mathbb{E}_{q(z|x)}[\log p(x|z)]}_{\text{\color{greencol}重构项}}
     - \underbrace{D_{\text{KL}}\!\big(q(z|x)\,\|\,p(z)\big)}_{\text{\color{vqcol}KL 正则项}}$};
\end{tikzpicture}
\end{frame}

% ── Slide 4: VQ-VAE 架构 ────────────────────────────
\begin{frame}{VQ-VAE 架构}
\centering
\begin{tikzpicture}[node distance=1cm, scale=0.88, every node/.style={scale=0.88}]
  % 主流程
  \node[vqb, minimum width=1.6cm] (input) {Input $x$};
  \node[vqb, right=1cm of input, minimum width=2cm] (enc) {Encoder};
  \node[vqb, right=1cm of enc, minimum width=1.6cm, fill=vqcol!75] (ze) {$z_e(x)$};
  \node[orb, right=1cm of ze, minimum width=2cm,
        minimum height=1.2cm] (quant) {\begin{tabular}{c}\small 量化\\[-2pt] \small Nearest\\[-2pt] \small Neighbor\end{tabular}};
  \node[vqb, right=1cm of quant, minimum width=1.6cm, fill=vqcol!90] (zq) {$z_q$};
  \node[decb, right=1cm of zq, minimum width=2cm] (dec) {Decoder};

  % 箭头
  \draw[arr] (input) -- (enc);
  \draw[arr=vqcol] (enc) -- (ze);
  \draw[arr] (ze) -- (quant);
  \draw[arr] (quant) -- (zq);
  \draw[arr] (zq) -- (dec);

  % 码本
  \node[above=1.2cm of quant, draw=orangecol, thick, rounded corners=5pt,
        fill=orangecol!8, minimum width=3cm, minimum height=2cm] (cb) {};
  \node[above=0.05cm of cb, font=\bfseries\color{orangecol}] {Codebook};
  % 码本向量可视化
  \foreach \i/\j/\hi in {-0.8/0.3/0, -0.2/0.3/0, 0.4/0.3/1,
                          -0.8/-0.2/0, -0.2/-0.2/0, 0.4/-0.2/0,
                          -0.8/-0.7/0, -0.2/-0.7/0, 0.4/-0.7/0}{
    \ifnum\hi=1
      \fill[vqcol] ([xshift=\i cm, yshift=\j cm]cb.center) rectangle
                   ++(0.4, 0.35);
    \else
      \fill[orangecol!40] ([xshift=\i cm, yshift=\j cm]cb.center) rectangle
                          ++(0.4, 0.35);
    \fi
  }
  \node[right=0.2cm of cb, font=\bfseries\color{vqcol}] {$e_k$};
  \draw[arr=orangecol] (cb) -- (quant);

  % straight-through 标注
  \node[below=0.5cm of quant, font=\small\itshape\color{purplecol}] (st)
    {straight-through estimator};
  \draw[dasharr=purplecol] ([yshift=-2pt]zq.south) -- ++(0,-0.3) -|
    ([yshift=-2pt]ze.south);
  \node[below=0.25cm of st, font=\scriptsize\color{purplecol}]
    {$\nabla_{z_e} \approx \nabla_{z_q}$};

  % 离散空间标注
  \node[below=2.2cm of quant, fill=lightred, rounded corners=4pt,
        inner sep=5pt, font=\bfseries\color{vqcol}] (dlabel)
    {离散潜在空间 (码本索引 $k \in \{1,\dots,K\}$)};

  % 损失
  \node[below=3.3cm of quant,
        draw=vqcol, rounded corners=5pt, fill=vqcol!5,
        inner sep=7pt, font=\small, text=darkcol] (loss)
    {$\displaystyle \mathcal{L}_{\text{VQ}}
     = \underbrace{\|x - \hat{x}\|^2}_{\text{\color{greencol}重构}}
     + \underbrace{\|\operatorname{sg}[z_e] - e_k\|^2}_{\text{\color{orangecol}码本更新}}
     + \beta\underbrace{\|z_e - \operatorname{sg}[e_k]\|^2}_{\text{\color{purplecol}commitment}}$};
\end{tikzpicture}
\end{frame}

% ── Slide 5: 潜在空间可视化 ─────────────────────────
\begin{frame}{潜在空间对比：连续 vs 离散}
\begin{columns}[T]
\begin{column}{0.48\textwidth}
\centering
{\large\bfseries\color{vaecol} VAE: 连续高斯空间}\\[8pt]
\begin{tikzpicture}[scale=0.75]
  % 坐标轴
  \draw[-{Stealth}, thick, darkcol!50] (-3.5,0) -- (3.5,0) node[right]{\small $z_1$};
  \draw[-{Stealth}, thick, darkcol!50] (0,-3.5) -- (0,3.5) node[above]{\small $z_2$};
  % 等高线圆
  \foreach \r/\op in {1/0.25, 2/0.15, 3/0.08}{
    \draw[vaecol, thick, opacity=\op, fill=vaecol, fill opacity=\op]
      (0,0) circle (\r);
  }
  % 散点 (手动放置代表性采样点)
  \foreach \x/\y in {0.3/0.5, -0.7/0.2, 0.1/-0.4, -0.3/-0.6,
                     1.2/0.8, -1.0/1.1, 0.8/-1.3, -1.5/-0.5,
                     1.8/0.2, -0.4/1.7, 0.5/-1.8, -1.2/-1.4,
                     2.1/-0.9, -1.8/1.3, 0.9/2.0, -2.2/0.1,
                     0.2/0.1, -0.1/-0.2, 0.6/0.3, -0.5/0.7}{
    \fill[vaecol!70] (\x, \y) circle (2.5pt);
  }
  % 标注
  \node[below, font=\small\itshape\color{vaecol}] at (0, -3.8)
    {$z = \mu + \sigma \cdot \varepsilon$，任意点可采样};
\end{tikzpicture}
\end{column}
\begin{column}{0.48\textwidth}
\centering
{\large\bfseries\color{vqcol} VQ-VAE: 离散码本空间}\\[8pt]
\begin{tikzpicture}[scale=0.75]
  % 坐标轴
  \draw[-{Stealth}, thick, darkcol!50] (-3.5,0) -- (3.5,0) node[right]{\small $z_1$};
  \draw[-{Stealth}, thick, darkcol!50] (0,-3.5) -- (0,3.5) node[above]{\small $z_2$};
  % Voronoi 边界 (近似)
  \foreach \x in {-2, 0, 2}{
    \draw[darkcol!15, thin] (\x-1, -3.3) -- (\x-1, 3.3);
  }
  \foreach \y in {-2, 0, 2}{
    \draw[darkcol!15, thin] (-3.3, \y-1) -- (3.3, \y-1);
  }
  % 码本点
  \foreach \x/\y in {-2/-2, -2/0, -2/2, 0/-2, 0/0, 0/2, 2/-2, 2/0, 2/2}{
    \fill[vqcol] (\x, \y) ++(-0.18,-0.18) rectangle ++(0.36,0.36);
  }
  % 编码器输出被 snap
  \foreach \ex/\ey/\qx/\qy in {
    -1.3/0.4/-2/0, 0.7/-1.5/0/-2, 1.8/0.3/2/0,
    -0.5/2.3/0/2, -2.4/-1.2/-2/-2, 1.1/1.6/2/2,
    -1.7/1.8/-2/2, 0.3/0.4/0/0, 2.3/-1.8/2/-2}{
    \fill[orangecol!60] (\ex, \ey) circle (2pt);
    \draw[darkcol!25, thin, -{Stealth[length=3pt]}] (\ex, \ey) -- (\qx, \qy);
  }
  % 标注
  \node[below, font=\small\itshape\color{vqcol}] at (0, -3.8)
    {$z_q = \arg\min_k \|z_e - e_k\|$，只有码本点};
\end{tikzpicture}
\end{column}
\end{columns}

\vspace{8pt}
\centering
\begin{tikzpicture}
  \fill[vqcol] (0,0) ++(-0.12,-0.12) rectangle ++(0.24,0.24);
  \node[right, font=\small] at (0.2, 0) {码本向量 $e_k$};
  \fill[orangecol!60] (3, 0) circle (2.5pt);
  \node[right, font=\small] at (3.2, 0) {编码器输出 $z_e$ (被 snap 到最近码本)};
\end{tikzpicture}
\end{frame}

% ── Slide 6: 梯度传播 ───────────────────────────────
\begin{frame}{梯度传播机制对比}
\begin{columns}[T]
\begin{column}{0.48\textwidth}
\centering
\vaebox{\large\bfseries\color{vaecol} VAE: 重参数化技巧}\\[10pt]
\begin{tikzpicture}[scale=0.85, every node/.style={scale=0.85}]
  \node[vaeb, minimum width=2cm] (enc) {Encoder};
  \node[vaeb, above right=0.4cm and 1cm of enc, minimum width=1.2cm, fill=vaecol!75] (mu) {$\mu$};
  \node[vaeb, below right=0.4cm and 1cm of enc, minimum width=1.2cm, fill=vaecol!75] (sig) {$\sigma$};
  \node[genb, right=2.2cm of enc, minimum width=2cm] (z) {$z$};
  \node[above=0.4cm of z, font=\small\itshape\color{greencol}] (eps) {$\varepsilon \sim \mathcal{N}(0,I)$};

  \draw[arr=vaecol] (enc) -- (mu);
  \draw[arr=vaecol] (enc) -- (sig);
  \draw[arr=greencol] (mu) -- (z);
  \draw[arr=greencol] (sig) -- (z);
  \draw[arr=greencol] (eps) -- (z);

  % 梯度反传
  \draw[dasharr=vaecol, line width=1.5pt]
    ([yshift=-8pt]z.west) -- ([yshift=-8pt]enc.east)
    node[midway, below, font=\small\color{vaecol}] {$\nabla$ 可反传};
\end{tikzpicture}

\vspace{10pt}
\begin{tikzpicture}
  \node[infobox=vaecol, text width=5.8cm] {
    \textbf{核心思想：}将随机性移到外部\\[3pt]
    $z = \mu + \sigma \cdot \varepsilon$ 中 $\varepsilon$ 是常量\\[3pt]
    $\Rightarrow$ $\frac{\partial z}{\partial \mu} = 1$,\;
    $\frac{\partial z}{\partial \sigma} = \varepsilon$\\[3pt]
    \color{greencol}\textbf{梯度畅通无阻 \checkmark}
  };
\end{tikzpicture}
\end{column}

\begin{column}{0.48\textwidth}
\centering
\vqbox{\large\bfseries\color{vqcol} VQ-VAE: Straight-Through}\\[10pt]
\begin{tikzpicture}[scale=0.85, every node/.style={scale=0.85}]
  \node[vqb, minimum width=1.5cm, fill=vqcol!75] (ze) {$z_e$};
  \node[orb, right=1cm of ze, minimum width=2cm] (q) {\small argmin};
  \node[vqb, right=1cm of q, minimum width=1.5cm] (zq) {$z_q$};

  \draw[arr] (ze) -- (q);
  \draw[arr] (q) -- (zq);

  % 前向
  \node[above=0.8cm of q, font=\small\bfseries\color{orangecol}] {前向: 走量化};

  % 反向 (straight-through)
  \draw[dasharr=purplecol, line width=1.5pt]
    ([yshift=-10pt]zq.west) -- ([yshift=-10pt]ze.east)
    node[midway, below, font=\small\color{purplecol}] {$\nabla$ 直接复制};
  \node[below=1cm of q, font=\small\bfseries\color{purplecol}] {反向: 跳过 argmin};
\end{tikzpicture}

\vspace{10pt}
\begin{tikzpicture}
  \node[infobox=vqcol, text width=5.8cm] {
    \textbf{核心思想：}前向用量化，反向绕过去\\[3pt]
    argmin 不可微 $\Rightarrow$ 梯度无法流过\\[3pt]
    Straight-through: $\nabla_{z_e} \approx \nabla_{z_q}$\\[3pt]
    \color{purplecol}\textbf{近似方法，工程上很有效 \checkmark}
  };
\end{tikzpicture}
\end{column}
\end{columns}
\end{frame}

% ── Slide 7: 损失函数详解 ───────────────────────────
\begin{frame}{损失函数详解}

% VAE Loss
\begin{center}
{\large\bfseries\color{vaecol} VAE Loss}
\end{center}
\vspace{-4pt}
\begin{tikzpicture}
  \node[draw=vaecol, thick, rounded corners=6pt, fill=lightblue,
        inner sep=10pt, text width=14cm, align=center] {
    $\displaystyle \mathcal{L}_{\text{VAE}}
     = \underbrace{\mathbb{E}_{q(z|x)}[\log p(x|z)]}_{\text{\color{greencol}\bfseries 重构损失：让输出接近输入}}
     - \underbrace{D_{\text{KL}}\!\big(q(z|x)\,\|\,p(z)\big)}_{\text{\color{vqcol}\bfseries KL 正则：让后验接近先验 $\mathcal{N}(0,I)$}}$
  };
\end{tikzpicture}

\vspace{10pt}

% VQ-VAE Loss
\begin{center}
{\large\bfseries\color{vqcol} VQ-VAE Loss}
\end{center}
\vspace{-4pt}
\begin{tikzpicture}
  \node[draw=vqcol, thick, rounded corners=6pt, fill=lightred,
        inner sep=10pt, text width=14cm, align=center] (vqloss) {
    $\displaystyle \mathcal{L}_{\text{VQ}}
     = \underbrace{\|x - \hat{x}\|^2}_{\text{\color{greencol}\bfseries 重构}}
     + \underbrace{\|\operatorname{sg}[z_e] - e_k\|^2}_{\text{\color{orangecol}\bfseries 码本更新}}
     + \beta\underbrace{\|z_e - \operatorname{sg}[e_k]\|^2}_{\text{\color{purplecol}\bfseries commitment}}$
  };
\end{tikzpicture}

\vspace{6pt}
\begin{columns}[T]
\begin{column}{0.32\textwidth}
\centering
\begin{tikzpicture}
\node[infobox=greencol, text width=4cm] {
  \textbf{\color{greencol}重构损失}\\
  解码器输出尽量\\接近原始输入
};
\end{tikzpicture}
\end{column}
\begin{column}{0.32\textwidth}
\centering
\begin{tikzpicture}
\node[infobox=orangecol, text width=4cm] {
  \textbf{\color{orangecol}码本损失}\\
  让码本向量 $e_k$ 靠近\\编码器输出（只更新码本）
};
\end{tikzpicture}
\end{column}
\begin{column}{0.32\textwidth}
\centering
\begin{tikzpicture}
\node[infobox=purplecol, text width=4cm] {
  \textbf{\color{purplecol}Commitment 损失}\\
  让编码器输出靠近\\码本向量（只更新编码器）
};
\end{tikzpicture}
\end{column}
\end{columns}

\vspace{4pt}
\centering
{\small\itshape\color{darkcol!60}
$\operatorname{sg}[\cdot]$ = stop-gradient：阻止梯度流过该项，分别更新码本与编码器}
\end{frame}

% ── Slide 8: 优缺点对比 ─────────────────────────────
\begin{frame}{优缺点对比}
\begin{columns}[T]
\begin{column}{0.47\textwidth}
\centering
{\Large\bfseries\color{vaecol} VAE}\\[10pt]

\begin{tikzpicture}
\node[draw=greencol, thick, rounded corners=5pt, fill=greencol!5,
      text width=6cm, inner sep=8pt, align=left] {
  \textbf{\color{greencol} 优点}\\[4pt]
  $\checkmark$ 理论优美，ELBO 有清晰概率解释\\[2pt]
  $\checkmark$ 可直接从先验 $p(z)$ 采样生成\\[2pt]
  $\checkmark$ 潜在空间平滑，支持插值\\[2pt]
  $\checkmark$ 训练简单稳定
};
\end{tikzpicture}

\vspace{8pt}
\begin{tikzpicture}
\node[draw=vqcol, thick, rounded corners=5pt, fill=vqcol!5,
      text width=6cm, inner sep=8pt, align=left] {
  \textbf{\color{vqcol} 缺点}\\[4pt]
  $\times$ 生成图像偏模糊（高斯假设）\\[2pt]
  $\times$ 容易 posterior collapse\\[2pt]
  $\times$ KL 与重构的平衡困难
};
\end{tikzpicture}
\end{column}

\begin{column}{0.47\textwidth}
\centering
{\Large\bfseries\color{vqcol} VQ-VAE}\\[10pt]

\begin{tikzpicture}
\node[draw=greencol, thick, rounded corners=5pt, fill=greencol!5,
      text width=6cm, inner sep=8pt, align=left] {
  \textbf{\color{greencol} 优点}\\[4pt]
  $\checkmark$ 生成质量更高更清晰\\[2pt]
  $\checkmark$ 避免 posterior collapse\\[2pt]
  $\checkmark$ 离散表示适配自回归模型\\[2pt]
  $\checkmark$ 可扩展至 VQ-VAE-2 多层级
};
\end{tikzpicture}

\vspace{8pt}
\begin{tikzpicture}
\node[draw=vqcol, thick, rounded corners=5pt, fill=vqcol!5,
      text width=6cm, inner sep=8pt, align=left] {
  \textbf{\color{vqcol} 缺点}\\[4pt]
  $\times$ 不能直接采样，需额外先验模型\\[2pt]
  $\times$ 码本利用率可能很低\\[2pt]
  $\times$ 训练更复杂（三项 loss 平衡）
};
\end{tikzpicture}
\end{column}
\end{columns}
\end{frame}

% ── Slide 9: 生成流程对比 ───────────────────────────
\begin{frame}{生成（推理）流程对比}

\vspace{6pt}
% VAE
{\large\bfseries\color{vaecol} VAE：一步到位}\\[6pt]
\centering
\begin{tikzpicture}[node distance=0.8cm, scale=0.9, every node/.style={scale=0.9}]
  \node[vaeb, minimum width=2.2cm] (s1) {$z \sim \mathcal{N}(0, I)$};
  \node[decb, right=1.5cm of s1, minimum width=2.2cm] (s2) {Decoder};
  \node[genb, right=1.5cm of s2, minimum width=2.2cm] (s3) {生成 $\hat{x}$};
  \draw[arr, line width=2pt] (s1) -- (s2);
  \draw[arr, line width=2pt] (s2) -- (s3);
  \node[below=0.3cm of s2, font=\small\itshape\color{vaecol}]
    {从先验采样 $\rightarrow$ 解码，一步完成};
\end{tikzpicture}

\vspace{16pt}
% VQ-VAE
\flushleft
{\large\bfseries\color{vqcol} VQ-VAE：两阶段}\\[6pt]
\centering
\begin{tikzpicture}[node distance=0.6cm, scale=0.9, every node/.style={scale=0.9}]
  % Stage 1 标注
  \node[font=\bfseries\color{darkcol}, anchor=west] (t1) at (-1, 1.5)
    {Stage 1: 训练 VQ-VAE};
  \node[font=\small\color{darkcol!60}, anchor=west] at (4, 1.5)
    {--- 学会编码 + 码本 + 解码};

  % Stage 2 标注
  \node[font=\bfseries\color{darkcol}, anchor=west] (t2) at (-1, 0.6)
    {Stage 2: 训练先验模型};
  \node[font=\small\color{darkcol!60}, anchor=west] at (4, 0.6)
    {--- 用自回归模型学码本索引分布};

  % 生成流程
  \node[purb, minimum width=2.2cm] (p1) at (0, -0.8) {PixelCNN};
  \node[right=0.3cm of p1, font=\small\color{darkcol}] (arr1) {};
  \node[orb, right=1cm of p1, minimum width=2.5cm] (p2) {\small 码本查找 $e_{k_i}$};
  \node[decb, right=1cm of p2, minimum width=2.2cm] (p3) {Decoder};
  \node[genb, right=1cm of p3, minimum width=2.2cm] (p4) {生成 $\hat{x}$};

  \draw[arr, line width=2pt] (p1) -- node[above, font=\scriptsize\color{darkcol}] {索引} (p2);
  \draw[arr, line width=2pt] (p2) -- node[above, font=\scriptsize\color{darkcol}] {$z_q$} (p3);
  \draw[arr, line width=2pt] (p3) -- (p4);

  \node[below=0.3cm of p2, font=\small\itshape\color{vqcol}]
    {先验模型采样索引 $\rightarrow$ 查码本 $\rightarrow$ 解码};
\end{tikzpicture}
\end{frame}

% ── Slide 10: 总结 ──────────────────────────────────
{
\setbeamercolor{background canvas}{bg=titlebg}
\begin{frame}[plain]
\vspace{0.8cm}
\begin{center}
{\Huge\bfseries\color{white} 总\quad 结}\\[16pt]
{\color{vqcol}\rule{4cm}{1.5pt}}
\end{center}
\vspace{6pt}

\begin{tikzpicture}
  \foreach \i/\title/\desc/\col in {
    1/{潜在空间}/{VAE 连续高斯\quad vs\quad VQ-VAE 离散码本}/vaecol,
    2/{核心机制}/{重参数化技巧\quad vs\quad Straight-Through Estimator}/greencol,
    3/{损失函数}/{KL 散度正则\quad vs\quad 码本损失 + Commitment 损失}/orangecol,
    4/{生成方式}/{直接先验采样\quad vs\quad 需额外自回归先验模型}/purplecol,
    5/{应用趋势}/{VQ-VAE 已成为 DALL-E、Stable Diffusion 等的基础组件}/vqcol%
  }{
    \pgfmathsetmacro{\ypos}{5.8 - (\i-1)*1.1}
    % 序号圆
    \fill[\col] (1, \ypos) circle (0.35);
    \node[font=\large\bfseries, text=white] at (1, \ypos) {\i};
    % 文字
    \node[anchor=west, font=\large\bfseries, text=white] at (1.6, \ypos+0.15) {\title};
    \node[anchor=west, font=\small, text=white!75] at (1.6, \ypos-0.25) {\desc};
  }
\end{tikzpicture}

\vspace{8pt}
\begin{center}
\begin{tikzpicture}
  \node[rounded corners=6pt, fill=vqcol!20, inner sep=10pt,
        font=\large\bfseries, text=white] {
    VAE: 连续 + KL \quad$|$\quad VQ-VAE: 离散 + 码本
    \quad$\longrightarrow$\quad 各有所长，互为补充
  };
\end{tikzpicture}
\end{center}
\end{frame}
}

\end{document}